\clearpage
\VerdierPhysicalPageAnchor{0190}%
\begingroup\small\setlength{\parskip}{0pt}\setlength{\emergencystretch}{1.5em}
\setlength{\abovedisplayskip}{4pt plus 2pt minus 2pt}%
\setlength{\belowdisplayskip}{4pt plus 2pt minus 2pt}%

\VerdierSemanticAnchor{ch3d-2-1-5}
\noindent\textbf{注意 2.1.5.}\quad
補題 (2.1.4) の仮定の下で、さらに
$X,Y\in\operatorname{Ob}\operatorname{comp}^{-}(\mathcal A)$ とする。
上の議論により、必要なら $X$ をホモトピー同値な複体で置き換えて、$m$ がエピ射で
あると仮定できる。上で与えた列 $X_n$ の構成から、帰納系 $X_n$ は定常で
あり（すなわち $X_n=Y$ かつ $u_n=\operatorname{id}$ が十分小さい $n$ に
対して成り立つ）、各 $n$ に対して
$X_n\in\operatorname{Ob}\operatorname{comp}^{-}(\mathcal A)$ である。

\VerdierSemanticAnchor{ch3d-2-2}
\subsection*{2.2. ある対象集合の考察。}

\VerdierSemanticAnchor{ch3d-2-2-1}
\noindent\textbf{2.2.1.}\quad
$M$ を $\mathcal A$ の対象の集合とする。次の諸性質を考える：

\VerdierSemanticAnchor{ch3d-2-2-1-e1}
\noindent $(E_1)$ 零対象は $M$ に属する。$M$ に属する二対象の直和は $M$ に
属する。

\VerdierSemanticAnchor{ch3d-2-2-1-e2}
\noindent $(E_2)$ 任意のエピ射 $X\xrightarrow{u}L$ に対して、ここで
$X$ は $\mathcal A$ の対象、$L$ は $M$ の元とし、エピ射
$v\colon L'\longrightarrow L$ で $L'\in M$ となり、かつ $u$ を経由する
ものが存在する。

\VerdierSemanticAnchor{ch3d-2-2-1-e3}
\noindent $(E_3)$ 任意の整数 $n$ と任意の完全列
\[
0\longrightarrow X\longrightarrow L_n\longrightarrow L_{n-1}
\longrightarrow\cdots\longrightarrow L_0\longrightarrow0,
\qquad L_i\in M,\quad 0\le i\le n,
\]
に対して、エピ射 $L'\longrightarrow X$ で $L'\in M$ となるものが存在する。

\VerdierSemanticAnchor{ch3d-2-2-1-e4}
\noindent $(E_4)$ 任意のエピ射 $u\colon L\longrightarrow L'$ に対して、ここで
$L,L'\in M$ とし、$u$ の核は $M$ に属する。

\VerdierSemanticAnchor{ch3d-2-2-1-e5}
\noindent $(E_5)$ 任意の拡大
$0\longrightarrow L\longrightarrow Z\longrightarrow L'\longrightarrow0$ に
対して、ここで $L,L'\in M$ とし、対象 $Z$ は $M$ に属する；$0$ は $M$ に
属する。

\VerdierSemanticAnchor{ch3d-2-2-1-e6}
\noindent $(E_6)$ 集合 $M$ は完全である。すなわち、$\mathcal A$ の短完全列の
三対象のうち二つが $M$ に属すれば第三の対象も同様である；$0$ は $M$ に
属する。

\VerdierSemanticAnchor{ch3d-2-2-1-implications}
含意 $(E_6)\Rightarrow(E_4)$ および $(E_5)$、
$(E_4)\Rightarrow(E_3)$、$(E_5)\Rightarrow(E_1)$ に注意する。両側無限完全列

\VerdierSemanticAnchor{ch3d-2-2-1-infinite-presentation}
\[
\cdots\longrightarrow L_n\longrightarrow L_{n-1}\longrightarrow\cdots
\longrightarrow L_0\longrightarrow X\longrightarrow0,
\qquad L_i\in M,\quad i\in\mathbb N,
\]
を\emph{$\infty$-$M$-表示}と呼ぶ。この表示の対象は $X$ である。$M$ を性質
$(E_1)$、$(E_2)$、$(E_3)$ をもつ対象の集合とする。$M^{\wedge}$ を次の
対象の集合とする：対象 $X$ が $M^{\wedge}$ に属するとは、
$\infty$-$M$-表示をもち、かつ任意の整数 $n$ と任意の完全列

\VerdierSemanticAnchor{ch3d-2-2-1-m-wedge}
\[
0\longrightarrow X'\longrightarrow L_n\longrightarrow L_{n-1}
\longrightarrow\cdots\longrightarrow L_0\longrightarrow X\longrightarrow0,
\qquad L_i\in M,\quad 0\le i\le n,
\]
に対して対象 $X'$ が $\infty$-$M$-表示をもつことをいう。
\endgroup
